% Generated draft: requires mathematical and source review.
\documentclass[11pt]{article}
\usepackage[margin=25mm]{geometry}
\usepackage{fontspec}
\setmainfont{Latin Modern Roman}
\usepackage{amsmath,amssymb,mathtools,unicode-math}
\setmathfont{Latin Modern Math}
\usepackage{enumitem,needspace,xurl,hyperref}
\setlength{\parindent}{0pt}
\setlength{\parskip}{5pt}
\setlength{\emergencystretch}{3em}
\newcommand{\N}{\mathbb{N}}
\newcommand{\Z}{\mathbb{Z}}
\newcommand{\Q}{\mathbb{Q}}
\newcommand{\R}{\mathbb{R}}
\newcommand{\C}{\mathbb{C}}
\newcommand{\F}{\mathbb{F}}
\newcommand{\A}{\mathbb{A}}
\newcommand{\PP}{\mathbb{P}}
\newcommand{\E}{\mathbb{E}}
\newcommand{\eps}{\varepsilon}
\newcommand{\dd}{\,\mathrm{d}}
\begin{document}
\section*{UnsolvedMath ID: 700721}
\textbf{Source problem code:} catalogue entry
\subsection{Definitions and mathematical statement}
For a finite group \(G\), let \(\pi(G)\) denote the set of prime divisors of \(|G|\). The prime graph of \(G\) has vertex set \(\pi(G)\); two distinct vertices \(p,q\) are adjacent if \(G\) contains an element of order \(pq\). Its disconnectedness means that there is a partition
\[
\pi(G)=\pi_1\sqcup\pi_2,\qquad \pi_1\ne\varnothing,\quad \pi_2\ne\varnothing,
\]
such that no element of \(G\) has order \(pq\) for any \(p\in\pi_1\) and \(q\in\pi_2\).

Let \(\Z G\) be the integral group ring, regarded as its regular left module, and interpret \(\Z\) as the trivial left \(\Z G\)-module, on which every element of \(G\) acts as the identity. A module is projective if it is isomorphic to a direct summand of a free module.

\textbf{Requested statement and proof requirement.} For every finite group \(G\) whose prime graph is disconnected, there exist left \(\Z G\)-modules \(A\) and \(B\) such that
\[
\Z\oplus\Z G\cong A\oplus B
\]
as \(\Z G\)-modules, with neither \(A\) nor \(B\) projective. Find a proof independent of the classification of finite simple groups (CFSG).

The supplied source attributes a proof using CFSG to P. A. Linnell. Its displayed module decomposition has malformed LaTeX; the isomorphism above is the natural reading of that display. The trivial action on \(\Z\) and the choice of left modules are interpretive conventions, since the supplied statement does not explicitly specify these details.
\subsection{Short English statement}
K. Gruenberg asks for a proof independent of the classification of finite simple groups that, whenever a finite group \(G\) has disconnected prime graph, \(\Z\oplus\Z G\) decomposes as a direct sum of two nonprojective \(\Z G\)-modules. Here \(\Z\) is interpreted as the trivial module. The source attributes a proof using the classification to P. A. Linnell.
\subsection{Sources}
\begin{itemize}
\item Supplied local UnsolvedMath registry record, permanent ID \(700721\), exact source code \texttt{catalogue entry}, titled \emph{Kourovka Notebook Problem 11.28}. Its status is supplied as \texttt{open}; this label is not an independently verified assessment.
\item \emph{The Kourovka Notebook (Unsolved Problems in Group Theory)}, original problem 11.28, 11th issue (1990), proposed by K. Gruenberg. The supplied background identifies the source edition as the 21st issue edition (2026), arXiv:\texttt{1401.0300v48}. The supplied collection description names E. I. Khukhro and V. D. Mazurov as editors.
\item P. A. Linnell, \emph{Proc. London Math. Soc.} \textbf{47}, no. 1 (1983), 83--127. The supplied statement attributes the decomposition theorem, proved using CFSG, to this reference; no article title is supplied.
\end{itemize}
\end{document}
